\subsection*{Drill: Running Away (Negative)}
\label{drill:power_jam/running_away}

This drill requires a three wall, a jammer at at least one offence.
Optionally this may be extended to a four wall, and up to four offences.  
This drill may be paired with a Positive power-jam drill for the offensive skaters. 


Starting with the four-wall 8-10 feet down track from the offence with the jammer another 5 to 10 feet behind, the goal is to maintain pack while blocking the jammer, but keep the offence away from the wall.

The goal of the offence is to get their jammer out.
This may be through holding position to try to break pack, taking a goat and holding pack, or attempting a regular offence. 

The defence in this first drill should focus on running away while holding the inside line.
If the defence has the luxury of a fourth blocker, they may use one blocker as a bridge whose additional job is to communicate when the wall should move.
The engagement zone extends 20ft from the bridge, while the bridge may hold up to 10 feet ahead of the foremost offence.
For safer ranges, the bridge should hold at around 8 feet from the offence, while the wall should be 8-15 feet from the bridge.  


\begin{itemize}
\item If the offence runs up, the bridge should communicate and the wall move up to maintain distance from the jammer.
\item If the offence stops then bridge should hold position. The defensive wall should then seek to recycle the jammer back to the bridge and reset. 
\item Goats should use appropriate tactics (see a goating drill for that).
\end{itemize}



